% sec-A083831.tex -- k = 4 section of the unified kp+1
% paper. Standalone fragment, \input by
% A083833-UNIFIE-paper.tex.
% Environments and notation inherited from sec-carry.tex; the
% guards below are no-ops once defined there.
\providecommand{\IsPal}{\mathrm{IsPal}}
\makeatletter
\@ifundefined{theorem}{\newtheorem{theorem}{Theorem}}{}
\@ifundefined{corollary}{\newtheorem{corollary}[theorem]{Corollary}}{}
\makeatother

\section{Multiplier $k = 4$: A083831 and its image A083832}
\label{sec:k4}

Entry A083831~\cite{oeis} (A.~Murthy and M.~Srikanth,
May 2003; corrected and extended by R.~Zumkeller and
R.~Chandler) lists the palindromes $p$ such that
$q = 4p+1$ is again a palindrome:
$1, 2, 8, 88, 131, 141, 232, 242, 888, 8888, \ldots$
Commenting on the entry (April 4, 2018), Robert Israel
recorded the conjecture quoted in
Section~\ref{sec:intro}. We settle it for every length,
in both directions, together with the digit-alphabet
bound on $p$ and, as a corollary of the same carry
analysis, the alphabet of the companion sequence.

\begin{theorem}[Israel's conjecture on A083831, full form]
\label{thm:k4-even}
Let $m \ge 1$ be even and $\IsPal(p,m)$. Then $4p+1$ is
a palindrome if and only if $a_j = 8$ for every
$j < m$: the even-length members of A083831 are exactly
the repdigits $\underbrace{88\ldots8}_{m}$.
\end{theorem}

\begin{theorem}[digit alphabet of the source]
\label{thm:k4-dig}
Under the same hypotheses (any $m \ge 1$), every digit
of $p$ satisfies $a_j \le 4$ or $a_j = 8$.
\end{theorem}

\begin{corollary}[alphabet of the image, OEIS A083832]
\label{cor:k4-alphabet}
Every digit of $q = 4p+1$ lies in $\{1,2,3,5,6,9\}$.
\end{corollary}

\begin{proof}
Instantiating the shared carry chain at $k = 4$ gives
$c_j = \bigl\lfloor\bigl(4(p \bmod 10^j) +
1\bigr)/10^j\bigr\rfloor$, with $c_0 = 1$ and
$c_j \le 3$.
\begin{itemize}
\item \emph{Same length} ($m' = m$):
Lemma~\ref{lem:same} forces $m$ odd,
$a_0 \in \{1,2\}$, $a_j \in \{0,1,2\}$ at even positions
and $a_j \in \{3,4\}$ at odd positions. Output digits:
$4a_j + 1 \in \{1,5,9\}$ at even $j$ and
$4a_j \in \{12,16\}$ at odd $j$, i.e.\ digits
$\{2,6\}$.
\item \emph{Overflow} ($m' = m+1$). Since
$4a \bmod 10$ does not separate the carry, the direct
mirror law of the $k = 5$ case is unavailable; the
induction runs on the mirrored couple \emph{(A)}+\emph{(B)}
of Lemma~\ref{lem:overflow}, whose single fixed point is
$(a_{j+1}, c_{j+2}) = (8, 3)$: the chain is rigid in one
step. The head equation has the single solution
$a_0 = 8$ (checked digit by digit, the $m = 1$ boundary
$33 \bmod 10 = 3 = \lfloor 33/10 \rfloor$ included), so
$c_1 = 3$; propagating the fixed point, every digit is
$8$ and the carries freeze at $c_j = 3$:
$p = 88\ldots8$ and $q = 35\ldots53$ ($q = 33$ for
$m = 1$), digits in $\{3,5\}$.
\end{itemize}
For $m$ even the same-length case is impossible, so the
overflow case applies and the forward direction of
Theorem~\ref{thm:k4-even} follows; the converse
direction is the certified repdigit membership
(Table~\ref{tab:lean}). The union of the two digit
classes gives Theorem~\ref{thm:k4-dig}, and reading
$b_j$ in both cases gives
$\{1,2,5,6,9\} \cup \{3,5\} = \{1,2,3,5,6,9\}$,
Corollary~\ref{cor:k4-alphabet}.
\end{proof}

The published terms agree: the even-length members are
exactly $88$, $8888$, $888888$; the image family
$33, 353, 3553, 35553, \ldots$ is the overflow branch
and $525, 565, 929, 969, 52125, \ldots$ the same-length
one.
